\textbf{结论名称：} 直角三角形存在性——两线一圆法（已知两点找第三点构成直角）

\textbf{适用条件：} 平面内已知两定点 \(A\)、\(B\)，求作或求满足 \(\triangle ABC\) 为直角三角形的动点 \(C\)；常用于中考坐标系或网格中的动点存在性问题。

\textbf{变量说明：} 已知点 \(A(x_1,y_1)\)，\(B(x_2,y_2)\)，动点 \(C(x,y)\)。向量 \(\overrightarrow{AB}=(x_2-x_1,\;y_2-y_1)\)。

\textbf{核心公式：}
\[
\boxed{
\begin{aligned}
&\angle C=90^\circ \Leftrightarrow \overrightarrow{CA}\cdot\overrightarrow{CB}=0 \quad (\text{轨迹：以}AB\text{为直径的圆},C\neq A,B)\\
&\angle A=90^\circ \Leftrightarrow \overrightarrow{AB}\cdot\overrightarrow{AC}=0 \quad (\text{轨迹：过}A\text{且垂直于}AB\text{的直线},C\neq A)\\
&\angle B=90^\circ \Leftrightarrow \overrightarrow{BA}\cdot\overrightarrow{BC}=0 \quad (\text{轨迹：过}B\text{且垂直于}AB\text{的直线},C\neq B)
\end{aligned}
}
\]
即：动点 \(C\) 的轨迹是“两直线（过两端点的垂线）与一圆（以 \(AB\) 为直径的圆）”除去 \(A,B\) 两点。

\textbf{使用场景：} 中考坐标背景下“找点构成直角三角形”的存在性问题；尺规作图寻找直角顶点；与抛物线、双曲线综合的动点分类讨论。

\textbf{简单例子：} 已知 \(A(0,0)\)，\(B(4,0)\)，求满足 \(\triangle ABC\) 为直角三角形且 \(y_C>0\) 的点 \(C\)。
\begin{itemize}
\item 若 \(\angle C=90^\circ\)：以 \(AB\) 为直径的圆 \((x-2)^2+y^2=4\)，取 \(x=1\) 得 \(y=\sqrt{3}\)，故 \(C_1(1,\sqrt{3})\)。
\item 若 \(\angle A=90^\circ\)：过 \(A\) 的垂线为 \(x=0\)，取 \(C_2(0,3)\)。
\item 若 \(\angle B=90^\circ\)：过 \(B\) 的垂线为 \(x=4\)，取 \(C_3(4,2)\)。
\end{itemize}
三种情况对应“两线一圆”上分别取点，均满足要求。